% translator_note.tex — 译者说明（放在目录之后、正文之前）
% 4–6 段中文；体例说明见 STYLE_GUIDE.md
\section*{译者说明}
\addcontentsline{toc}{section}{译者说明}

一、底本。本译注本依据 Daniel Berwick-Evans（University of Illinois Urbana-Champaign 数学系）的综述 \textit{Elliptic Cohomology and Quantum Field Theory}（arXiv:2408.07693v1 [math.AT]，2024 年 8 月 14 日，34 页）全文译出。原文是一篇导引性的综述章节（原文自称 chapter），目标是把「椭圆上同调的上闭链即二维超对称场论」这一猜想（原文猜想 1.1）拆解为可操作的分析数据与配边数据，并用大量例子说明其来源与后果。译文保持原文的章节结构与论证顺序，不作改写、删节或补充；原文的段落划分、脚注位置与文献引用标记一仍其旧。

二、编号。原文全文的编号公式构成单一序列 (1)–(56)，本译本逐条保留原编号，以 \verb|\tag| 手工排出，与原文严格一一对应。原文的编号环境（猜想、定义、注、例、命题、定理、推论、假设）在同一节内共享计数器，本译本同样按原文号手工排出，不重排、不跳号；原文中带括号说明的条目（如「定义 1.5（[ST11, Definition 5.2] 之略述）」），译文亦保留该说明。正文引用一律写作「式 (12)」「猜想 1.1」「定义 2.13」这类形式，不使用 \verb|\eqref|，以便与原文的纸质阅读顺序一致。原文公式编号排在页面左侧边栏，译本沿用中文数学书的通行体例排在右侧。

三、图与表。原文的 8 幅插图（Figure 1–8）均不复刻，只在原位置保留「图从略」标记与图注译文，正文中「见图 3」一类指引照译。交换图与显示公式以 LaTeX 重排：能写成配边三角形、自然变换方块的一律按原样排出，2-胞之类无法用 CD 环境表达的图形用简化的 TikZ 图重排，其顶点、边与标签与原文一一对应。原文正文提到的 Table 2 在原文中并无对应实体表格，译文照译并在源码批注中说明。

四、批注。源码中以 \% [Tutor Note: …] 形式保留了 Tutor Note 批注，内容限于四类：可验证的易错点、原文省略论证的补充、原文笔误、记号说明。批注不进入排版成品，只随源码交付。凡原文确有的笔误，译文一律照原文排印，另在批注中指出正确形式，并在工程文件目录下的 ERRATA.md 中登记；可疑而不能确定的写法只加批注、不改动原文，以免以译者之见覆盖作者原意。原文的脚注共 9 处，译文用 LaTeX 自动编号，个别处与原文序号不同者已在批注中注明；原文首页不编号的作者单位与 NSF 资助说明（DMS-2205835、DMS-2340239）移至封面。

五、术语。术语以工程文件目录下的 glossary.tsv 冻结表为准（共 151 条），同一英文术语全文用同一中文译名。人名、理论名、软件名与文献标签不音译：Witten、Stolz、Teichner、Segal、Hopkins、Lurie、Atiyah–Singer、Quillen、Ochanine、TMF、K-theory、Spin、string、Ramond、Neveu–Schwarz、Wick 等保留拉丁字母。表外的新术语在首次出现处写作「中文（English）」，并已回填该表。数学表述以原文页面图像为唯一依据（原文 PDF 的文字层会把分式压成一行、把矩阵拆成单字符），每条编号公式均经放大核对，不凭文字层转写。

六、致谢与免责。本译本仅供学习与研究之用，著作权归原作者 Daniel Berwick-Evans 所有；译文中的任何错误与不当表述均由译者负责，欢迎指正。感谢原作者写下这篇清晰而慷慨的综述，也感谢 Stolz 与 Teichner 的框架（原文 [ST11]）为本译本的术语选择提供了依据。
